% !TEX program = lualatex
% Preview of odv Grotesk: the added Greek and maths glyphs in text, and equations in
% odv Math as odv.sty sets them. Build the fonts first (python3 scripts/build.py and
% python3 scripts/odvmath.py), then: lualatex preview.tex (or xelatex)
\documentclass[11pt]{article}
\usepackage[a4paper, left=25mm, right=45mm, top=25mm, bottom=25mm]{geometry}
\usepackage{xcolor}
\definecolor{ink}{HTML}{161514}
\definecolor{muted}{HTML}{6B655E}
\color{ink}
\usepackage{unicode-math}
\setmainfont{odvGrotesk}[
  Path = ../fonts/ttf/, Extension = .ttf,
  UprightFont = *-Regular, ItalicFont = *-Italic,
  BoldFont = *-Bold, BoldItalicFont = *-BoldItalic]
% Maths as odv.sty sets it: odv Math (New Computer Modern Math Book built at 112 %) at
% 100 %, \mathrm from New Computer Modern Book at 112 % (Latin Modern if not installed).
% \text{} stays odv Grotesk.
\setmathfont{odvMath-Regular.otf}[Path = ../fonts/otf/]
\IfFontExistsTF{NewCM10-Book.otf}
  {\setmathrm{NewCM10-Book.otf}[Scale = 1.12, BoldFont = NewCM10-Bold.otf]}
  {\setmathrm{lmroman10-regular.otf}[Scale = 1.12, BoldFont = lmroman10-bold.otf]}
\linespread{1.25}
\setlength{\parindent}{0pt}
\setlength{\parskip}{0.6\baselineskip}
\newcommand\caps[1]{{\addfontfeatures{LetterSpace=6}\footnotesize\color{muted}\MakeUppercase{#1}}}
\newfontfamily\R{odvGrotesk-Regular.ttf}[Path=../fonts/ttf/]
\newfontfamily\I{odvGrotesk-Italic.ttf}[Path=../fonts/ttf/]
\newfontfamily\M{odvGrotesk-Medium.ttf}[Path=../fonts/ttf/]
\newfontfamily\MI{odvGrotesk-MediumItalic.ttf}[Path=../fonts/ttf/]
\newfontfamily\B{odvGrotesk-Bold.ttf}[Path=../fonts/ttf/]
\newfontfamily\BI{odvGrotesk-BoldItalic.ttf}[Path=../fonts/ttf/]
\newfontfamily\SG{SchibstedGrotesk-Regular.ttf}[Path=../upstream/]
\pagestyle{empty}
\begin{document}

\caps{odv Grotesk · preview 2 · 13 November 2026}

{\fontsize{26}{30}\selectfont\M Greek and maths for Schibsted Grotesk}

\caps{added Greek}\\[2pt]
{\fontsize{20}{28}\selectfont
\R ΓΘΛΞΦΨ γδεϵζηθκλνξρσςτυφϕχψω\\
\I ΓΘΛΞΦΨ γδεϵζηθκλνξρσςτυφϕχψω\\
\M ΓΘΛΞΦΨ γδεϵζηθκλνξρσςτυφϕχψω\\
\MI ΓΘΛΞΦΨ γδεϵζηθκλνξρσςτυφϕχψω\\
\B ΓΘΛΞΦΨ γδεϵζηθκλνξρσςτυφϕχψω\\
\BI ΓΘΛΞΦΨ γδεϵζηθκλνξρσςτυφϕχψω}

\caps{added maths}\\[2pt]
{\fontsize{20}{28}\selectfont
\R ∇ ∈ ∉ ∀ ∃ ∓ ⋅ ⊥ ∩ ∪ ⊂ ⊃ ⊆ ⊇ ∧ ∨ ⇒ ⇔ ∝ ∥ ∠ △\\
\I ∇ ∈ ∉ ∀ ∃ ∓ ⋅ ⊥ ∩ ∪ ⊂ ⊃ ⊆ ⊇ ∧ ∨ ⇒ ⇔ ∝ ∥ ∠ △\\
\M ∇ ∈ ∉ ∀ ∃ ∓ ⋅ ⊥ ∩ ∪ ⊂ ⊃ ⊆ ⊇ ∧ ∨ ⇒ ⇔ ∝ ∥ ∠ △\\
\MI ∇ ∈ ∉ ∀ ∃ ∓ ⋅ ⊥ ∩ ∪ ⊂ ⊃ ⊆ ⊇ ∧ ∨ ⇒ ⇔ ∝ ∥ ∠ △\\
\B ∇ ∈ ∉ ∀ ∃ ∓ ⋅ ⊥ ∩ ∪ ⊂ ⊃ ⊆ ⊇ ∧ ∨ ⇒ ⇔ ∝ ∥ ∠ △\\
\BI ∇ ∈ ∉ ∀ ∃ ∓ ⋅ ⊥ ∩ ∪ ⊂ ⊃ ⊆ ⊇ ∧ ∨ ⇒ ⇔ ∝ ∥ ∠ △}

\caps{accents}\\[2pt]
{\fontsize{20}{28}\selectfont\R ά έ ή ί ό ύ ώ ϊ ϋ ΐ ΰ Ά Έ Ή Ί Ό Ύ Ώ Ϊ Ϋ · λέξη δέλτα ξένος ψυχή}

\caps{already in Schibsted Grotesk, for comparison}\\[2pt]
{\fontsize{20}{28}\selectfont\SG ΔΠΣΩ αβμπ ∂×± ok pv xn t}

\caps{next to the Latin letters they borrow from}\\[2pt]
{\fontsize{20}{28}\selectfont\R kκ pρ oσ vγ xχ nη tτ oδ uυ ∈ϵε}

\caps{in running text}

The derated thermal transmittance adds a correction ζ to the one-dimensional value U\textsubscript{0}.
Linear thermal bridges enter as l·Ψ and point bridges as n·χ, so the scope of this paper is the sum of the Ψ and χ terms; κ converts model units to metres.
The relative errors δ\textsubscript{phys} and δ\textsubscript{recon} stay below 2\,\% in every case, and ΔT\textsubscript{si,min} below 0.1\,K.
\emph{In italic: the flux density q follows Fourier's law with conductivity λ, so ζ, κ, χ and Ψ are all derived from ∇T.}

\caps{in equations}
\begin{equation}
    U_d = \underbrace{\frac{1}{R_\mathrm{si} + \sum_i^n \frac{d_i}{\lambda_{2,i}}+ R_\mathrm{se}}}_{U_0} \cdot \overbrace{\left( 1+ \zeta \right)}^{\mathrm{derating}}
\end{equation}
\begin{equation}
    U_{d} = U_0 + \Delta U_f + \Delta U_g + \Delta U_r + \sum_j^n l_j \Psi_j + \sum_k^m n_k \chi_k
\end{equation}
\begin{equation}
    \vec{q} = -\lambda \nabla T, \qquad \vec{q}(x,y,z) \in \mathbb{R}^3
\end{equation}
\begin{equation}
    \zeta = \frac{\iint_A \vec{q} \cdot \vec{n} \, \mathrm{d}A}{U_0 \cdot A \cdot \Delta T} - 1
\end{equation}
\begin{equation}
    U_{\mathrm{eff}} = \kappa\,\frac{V_{\mathrm{solid}}}{s\,A}, \qquad
    \delta_\mathrm{phys} = \left( \frac{Q_\mathrm{FEM}}{Q_\mathrm{ref}} - 1 \right) \cdot 100\,\%,
    \qquad \mathbf{n} = \frac{\mathbf{t}_u \times \mathbf{t}_v}{\|\mathbf{t}_u \times \mathbf{t}_v\|}
\end{equation}

\end{document}
